\section{Lecture 22 (November 27th)}
\begin{rmk}
Last class, we have learned the fundamental theorem of Galois theory I. Today, we will look at some examples and the concept of group action.
\end{rmk}
\vspace{2ex}
\begin{thm}
(15.31) For a finite Galois extension $F/K$, there exists a correspondence between the Galois group (the group of automorphisms) and intermediate fields. Explicitly, the map
\[\{K\subset E\subset F\}\leftrightarrow \{H\triangleleft \mathrm{Gal}(F/K)\}\]
is given by
\[E\mapsto \mathrm{Gal}(F/E)\]
and
\[H\mapsto F^{H}\]
\end{thm}
\vspace{2ex}
\begin{defi}
(15.22) Let $F/K$ be a finite Galois extension. Let $\alpha \in F$ be an element. The disinct $\sigma (\alpha )$ for $\sigma \in \mathrm{Gal}(F/K)$ are the Galois conjugates of $\alpha $. 
\end{defi}
\vspace{2ex}
\begin{prop}
(15.24) Let $F/K$ be a finite Galois extension. Let $\alpha \in F$ be an element. Let $\alpha_1,\ldots ,\alpha _{r}$ be the Galois conjugates of $\alpha $. Then
\[f(t)=\prod^{r}_{i=1}(t-\alpha _{i})\]
is the minimal polynomial of $\alpha $ over $K$.
\end{prop}
\vspace{2ex}
\begin{proof}
It suffices to show that $\mathrm{deg}\,f(t)=\mathrm{deg}\,\mathrm{irr}(\alpha ,K)$. Note that 
\[\begin{align*}
\mathrm{deg}(\mathrm{irr}(\alpha ,K))=[K(\alpha ):K]=\dfrac{[F:K]}{[F:K(\alpha )]}=\dfrac{|\mathrm{Gal}(F/K)|}{|\mathrm{Gal}(F:K(\alpha ))|}=[\mathrm{Gal}(F/K):\mathrm{Gal}(F/K(\alpha ))]
\end{align*}\]
Consider the action of $\mathrm{Gal}(F/K)$ on the set of the Galois conjugates of $\alpha $. Applying lemma (12.9), we see that
\[\mathrm{deg}\,f(t)=r=[\mathrm{Gal}(F/K):\mathrm{Gal}(F/K(\alpha ))]\]
\end{proof}
\vspace{2ex}
\begin{defi}
(11.9) (Group actions) Let $G$ be a group and $X$ be a set. An action of $G$ on $X$ is a function $\phi :G\times X\rightarrow X$ such that:
\begin{itemize}
\item[(i)] $\phi (1,x)=x$ for all $x\in X$.
\item[(ii)] $\phi (gh,x)=\phi (g,\phi (h,x))$ for all $g,h\in G$ and all $x\in X$.
\end{itemize}
\end{defi}
\vspace{2ex}
\begin{prop}
(11.20) Giving an action of $G$ on $X$ is equivalent to giving a group homomorphism 
\[G\rightarrow S_{x}\]
\end{prop}
\vspace{2ex}
\begin{proof}
Ours! 
\end{proof}
\vspace{2ex}
\begin{ex}
\begin{itemize}
\item[(i)] $G$ acts on itself by left multiplication $G\times \underline{G}\rightarrow G:(g,h)\mapsto gh$
\item[(ii)] $G$ acts on itself by conjugation $G\times \underline{G}\rightarrow \underline{G}:(g,h)\mapsto ghg^{-1}$ 
\end{itemize}
\end{ex}
\vspace{2ex}
\begin{defi}
(11.22) Let $X$ be a $G$-set. The orbit of an element $x\in X$ is the subset $\{g\cdot x\;|\;g\in G\}\subset X$.
\end{defi}
\vspace{2ex}
\begin{defi}
(12.8) The stabiliser of $x\in X$ is the subgroup
\[G_{x}=\{g\in G\;|\;g\cdot x=x\}\]
\end{defi}
\vspace{2ex}
\begin{prop}
Let $S_{n}$ act on $\{1,\ldots ,n\}$. Let $H\triangleleft S_{n}$. Then is $H=(S_{n})_{i}$ for some $i$?
\end{prop}
\vspace{2ex}
\begin{lem}
(12.9) There is a bijection between the orbit of $X$ and the set $G/G_{X}$ of left cosets of $G_{X}$. 
\end{lem}
\vspace{2ex}
\begin{proof}
Define a function $\phi :G_{X}\rightarrow G/G_{X}$. Check that it is well defined and that it is bijective. 
\end{proof}
\vspace{2ex}
\begin{prop}
Let $K$ be a field, and $f(t)\in K[t]$ be a separable polynomial. The following are equivalent:
\begin{itemize}
\item[(i)] $f(t)\in K[t]$ is irreducible
\item[(ii)] $\mathrm{Gal}(K(\alpha_1,\ldots ,\alpha _{n})/K)$ acts transistively on $\{\alpha _1,\ldots ,\alpha _{n}\}$ where $\alpha_1,\ldots ,\alpha _{n}$ are the zeros of $f(t)$ in $\bar{K}$. That is, for any given pair of elements $(\alpha _{i},\alpha _{j})$, there exists an element $\sigma \in \mathrm{Gal}(K(\alpha_1,\ldots ,\alpha _{n})/K)$ such that $\sigma (\alpha _{i})=\sigma(\alpha _{j})$.
\end{itemize}
\end{prop}
\vspace{2ex}
\begin{rmk}
Can we prove the irreducibility of $f(t)$ on page 44 using the above proposition?
\end{rmk}
\vspace{2ex}

\begin{ex}
(15.32) Let $F$ be the splitting field of $t^3-2\in {\bm Q}[t]$ which is equal to ${\bm Q}(\alpha_1,\alpha_2,\alpha_3)={\bm Q}(2^{1/3},\alpha_1\xi ,\alpha_2\xi )$.
\[\mathrm{Gal}(F/{\bm Q})\cong S_{3}\]
Due to the above, we can notice that
\begin{center}
\begin{tikzpicture}[scale=1.2, every node/.style={font=\small}]

% Bottom field
\node (Q) at (0,0) {$S_3$};

% First layer
\node (Qzeta) at (-2,2) {$\langle (1,2,3)\rangle $};
\node (Qcub2) at (6,2) {$\langle (3,1)\rangle $};
\node (Qs2) at (4,2) {$\langle (2,3)\rangle $};

% Top field
\node (top) at (0,4) {$1$};

% Edges to bottom
\draw (Q) -- node[left] {} (Qzeta);
\draw (Q) -- node[right] {} (Qcub2);
\draw (Q) -- node[right] {} (Qs2);

% Edges to top
\draw (Qzeta) -- node[left] {} (top);
\draw (Qcub2) -- node[left] {} (top);
\draw (Qs2) -- node[right] {} (top);

% Diagonal middle field
\node (Qs) at (2,2) {$\langle (1,2)\rangle $};

\draw (Q) -- node[right] {} (Qs);
\draw (Qs) -- node[right] {} (top);

\end{tikzpicture}
\end{center}
has a bijection with
\begin{center}
\begin{tikzpicture}[scale=1.2, every node/.style={font=\small}]

% Bottom field
\node (Q) at (0,0) {${\bm Q}$};

% First layer
\node (Qzeta) at (-2,2) {$F^{\langle (1,2,3)\rangle} $};
\node (Qcub2) at (6,2) {$F^{\langle (3,1)\rangle }$};
\node (Qs2) at (4,2) {$F^{\langle (2,3)\rangle} $};

% Top field
\node (top) at (0,4) {$F$};

% Edges to bottom
\draw (Q) -- node[left] {} (Qzeta);
\draw (Q) -- node[right] {} (Qcub2);
\draw (Q) -- node[right] {} (Qs2);

% Edges to top
\draw (Qzeta) -- node[left] {} (top);
\draw (Qcub2) -- node[left] {} (top);
\draw (Qs2) -- node[right] {} (top);

% Diagonal middle field
\node (Qs) at (2,2) {$F^{\langle (1,2)\rangle } $};

\draw (Q) -- node[right] {} (Qs);
\draw (Qs) -- node[right] {} (top);

\end{tikzpicture}
\end{center}
Letting us conclude that
\[[F:F^{H}]=|H|\]

\end{ex}
\vspace{2ex}

